\documentclass[11pt]{article}

\usepackage[T1]{fontenc}
\usepackage[utf8]{inputenc}
\usepackage{lmodern}
\usepackage{textcomp}
\usepackage[margin=1in]{geometry}
\usepackage{amsmath,amssymb,amsthm}
\usepackage{enumitem}
\usepackage{microtype}
\usepackage[hidelinks]{hyperref}

\setlength{\emergencystretch}{2em}
\numberwithin{equation}{section}

\newtheorem{theorem}{Theorem}[section]
\newtheorem{lemma}[theorem]{Lemma}
\newtheorem{proposition}[theorem]{Proposition}
\newtheorem{corollary}[theorem]{Corollary}
\theoremstyle{definition}
\newtheorem{definition}[theorem]{Definition}
\newtheorem{example}[theorem]{Example}
\theoremstyle{remark}
\newtheorem{remark}[theorem]{Remark}

\newcommand{\Fcal}{\mathcal F}
\newcommand{\Hcal}{\mathcal H}
\newcommand{\Inv}{\operatorname{Inv}}
\newcommand{\Sat}{\operatorname{Sat}}

\title{Forcing Graphs and Cyclic Completion Energy\\
in Two-Comparable Triple Systems}
\author{Bulau Robert Nicolae\\
\small Colegiul Na\textcommabelow{t}ional ``Mihai Viteazul'',
Ploie\textcommabelow{s}ti, Romania}
\date{}

\begin{document}
\maketitle

\begin{abstract}
Let $G(n)$ be the maximum size of a subset of $[n]^3$ in which every two
distinct triples are strictly ordered in at least two coordinates.  The best
known general upper bound is $o(n^2)$, and it remains open whether
$G(n)=O(n^{2-\delta})$ for an absolute $\delta>0$.  We study the repeated-coordinate
constraints behind this problem through its equivalent formulation by monotone
partial matrices.  Every pair of equal-label entries forces two cross holes.
We orient these holes canonically and obtain an edge-coloured bipartite forcing
graph: vertex degrees are forcing multiplicities, neighbourhoods are strict
product-order chains, colour classes are order-preserving matchings, and every
four-cycle is rainbow.  In the three coordinate representations, the second
completion energies satisfy the cyclic packing law
\[
 E_x+E_y+E_z\le \binom{|M|}{2}.
\]
An exact decomposition of the unused capacity records coordinate-degree
variance, forcing-degree variance, aligned pairs, and unsaturated inversions.
We also prove endpoint-exclusion inequalities for common fibres.  Finally, the
centres of all pressure grids through a fixed cell split into two signed sets,
each admitting two canonical encodings as a new two-comparable set.  This gives
pointwise bounds $L_0\le 2|M|$ and $L_1\le 2|M|^{3/2}$, with the exponent $3/2$
sharp, and yields a fixed power saving under a quantitative pressure-load
hypothesis.  These results impose simultaneous structural restrictions on a
hypothetical near-quadratic family but do not establish such a power saving.
\end{abstract}

\section{Introduction}

For $u,v\in[n]^3$, write $u\prec_2 v$ if $u$ is strictly smaller than $v$
in at least two coordinates.  A set $M\subseteq[n]^3$ is \emph{pairwise
two-comparable} if, for every distinct $u,v\in M$, either $u\prec_2v$ or
$v\prec_2u$.  We consider the extremal function
\[
 G(n)=\max\{|M|:M\subseteq[n]^3\text{ is pairwise two-comparable}\}.
\]
The projection bound $G(n)\le n^2$ is immediate, since two members of $M$
cannot agree in two coordinates.  The problem is nevertheless far from this
elementary bound.  The related two-increasing sequence problem was posed by
Loh~\cite{Loh}.  Gowers and Long proved, using quantitative triangle removal,
that
\[
 G(n)\le \frac{n^2}{\exp(\Omega(\log^*n))}=o(n^2),
\]
and gave constructions of size at least $n^{1.546}$ for arbitrarily large
$n$~\cite{GowersLong,Fox}.  The fixed-power question
\begin{equation}\label{eq:open-problem}
 G(n)=O(n^{2-\delta})\qquad\text{for some absolute }\delta>0
\end{equation}
remains open.  Equivalent formulations occur in monotone matrices and tripod
packing; see Stein~\cite{Stein}, Tiskin~\cite{Tiskin}, Östergård and
Pöllänen~\cite{OstergardPollanen}, and Aronov et al.~\cite{AronovEtAl}.
The distinction from a two-increasing \emph{sequence} is important: Gowers and
Long obtained a power saving for sequences by using acyclicity, which is absent
for general two-comparable sets.  Recent higher-dimensional results for the
related increasing-sequence and tournament problems likewise leave the
three-coordinate comparable-set problem separate~\cite{FoxSudakovWigderson}.

The present paper develops structural constraints on a two-comparable family.
Choose one coordinate as a label and the other two as row and column.  The
result is a partial matrix in which rows, columns, and equal-label positions
are all strictly increasing.  If a label occurs in both the row and the column
of an empty cell, then that label \emph{forces} the cell to remain empty.  The
main results organize these forcing events in four ways.

First, Section~\ref{sec:forcing-graph} constructs a canonical oriented forcing
graph.  It is a simple edge-coloured bipartite graph whose degrees and wedges
are the first and second completion counts.  Its neighbourhoods are ordered
chains, its colour classes are order-preserving matchings, and all four-cycles
and bicliques are rainbow.  Complete bipartite subgraphs of every fixed size
can nevertheless occur, so the structure is not captured by a fixed forbidden
subgraph.

Second, Section~\ref{sec:common-fibres} proves an endpoint-exclusion theorem for
any collection of common row fibres and derives a hierarchy of binomial-moment
inequalities.  Third, Sections~\ref{sec:cyclic-energy} and
\ref{sec:exact-slack} couple all three coordinate representations.  A
double-forcing event is put in bijection with a saturated inversion.  The
three inversion images are disjoint, giving
\begin{equation}\label{eq:intro-cyclic}
 E_x+E_y+E_z\le \binom{|M|}{2}.
\end{equation}
The associated exact slack identity retains every loss in the two variance
steps and in the cyclic packing step.  It gives an elementary density
calibration, an equality-stability statement, and a criterion showing that a
near-quadratic family must have almost maximal polynomial-scale forcing energy
in every coordinate representation.

Fourth, Section~\ref{sec:pressure} studies the overlap of the rectangular
pressure grids forced around holes.  For a fixed physical cell, the centres of
all pressure grids containing it split according to orientation.  Each signed
centre set has two canonical encodings as a pairwise two-comparable subset of
$[n]^3$.  A discrete Loomis--Whitney estimate then gives the pointwise load
bounds
\begin{equation}\label{eq:intro-load}
 L_0(z)\le 2|M|,\qquad L_1(z)\le 2|M|^{3/2}.
\end{equation}
The weighted exponent is sharp as a function of $|M|$.  Moreover, if the
maximum weighted pressure load is $O(n^{3-\eta})$, then
\[
 |M|=O(n^{2-\eta/6}+n).
\]
Thus a concrete overlap estimate would imply \eqref{eq:open-problem}; the
unconditional pointwise estimate identifies the overlap as a smaller copy of
the original two-comparable problem.

All results in the body of the paper are proved mathematically and do not
depend on computation.  Appendix~\ref{app:verification} describes an exhaustive
order-three verification used to check definitions, factors of two, and the
explicit constructions.  Section~\ref{sec:scope} records the precise limits of
the methods.  In particular, no absolute $\delta>0$ is obtained here.

\section{Monotone partial matrices and completion counts}
\label{sec:preliminaries}

We first give the matrix correspondence and fix notation.  An order-$n$
partial matrix $A$ has cells $[n]\times[n]$, some of which are filled by
labels from $[n]$.  Its sets of filled cells and holes are denoted by
$\Fcal(A)$ and $\Hcal(A)$.  It is \emph{monotone} if:
\begin{enumerate}[label=(\roman*)]
\item the filled entries in each row strictly increase from left to right;
\item the filled entries in each column strictly increase from top to bottom;
\item the cells carrying any fixed label form a strict chain in the
row--column product order.
\end{enumerate}
Here $(i,j)<(i',j')$ means $i<i'$ and $j<j'$.

Given a pairwise two-comparable set $M\subseteq[n]^3$, choose a coordinate
$q\in\{x,y,z\}$ as the label and use the remaining two coordinates, in their
cyclic order, as row and column.  Since two triples cannot agree in two
coordinates, this defines a partial matrix $A^{(q)}$.  A pair agreeing in the
row, column, or label coordinate gives, respectively, conditions (i), (ii), or
(iii).  Conversely, these three conditions guarantee two-comparability for
pairs agreeing in one coordinate; if all three coordinates differ, one of the
two triples is smaller in at least two coordinates by the majority principle.
We have therefore proved the following standard equivalence.

\begin{proposition}[Matrix correspondence]\label{prop:matrix-correspondence}
Pairwise two-comparable subsets of $[n]^3$ are in bijection with order-$n$
monotone partial matrices.  The bijection may use any coordinate as label.
\end{proposition}

Fix one representation $A$.  If a hole $h=(i,j)$ has a label $k$ occurring
both in row $i$ and in column $j$, then $k$ \emph{forces} $h$.  Write
\[
 K_h=\{k:k\text{ forces }h\},\qquad d(h)=|K_h|.
\]
The forcing multiplicity is zero for an unforced hole.  Let $t_k$ be the
number of occurrences of label $k$.  The first completion count is exact.

\begin{lemma}[First completion identity]\label{lem:first-completion}
For every monotone partial matrix,
\begin{equation}\label{eq:first-completion}
 \sum_{h\in\Hcal(A)}d(h)=\sum_{k=1}^n t_k(t_k-1).
\end{equation}
\end{lemma}

\begin{proof}
Order the $t_k$ occurrences of a fixed label $k$ as
$(r_1,c_1)<\cdots<(r_{t_k},c_{t_k})$.  For every $p<q$, the two cross cells
$(r_p,c_q)$ and $(r_q,c_p)$ are holes: filling either would contradict row and
column monotonicity.  Both are forced by $k$.  Conversely, a forcing event
$(h,k)$ uniquely determines the occurrence of $k$ in the row of $h$ and the
occurrence of $k$ in its column, hence one of these two cross cells.  Thus label
$k$ contributes $2\binom{t_k}{2}=t_k(t_k-1)$ events.  Summing over $k$ proves
the identity.
\end{proof}

The \emph{second completion energy} is
\begin{equation}\label{eq:second-energy}
 E(A)=\sum_{h\in\Hcal(A)}\binom{d(h)}2,
\end{equation}
where $\binom d2=0$ for $d\in\{0,1\}$.  It counts a hole together with an
unordered pair of labels that both force it.

\section{The oriented forcing graph}
\label{sec:forcing-graph}

We begin with the local orientation that makes the graph canonical.

\begin{lemma}[Common orientation]\label{lem:common-orientation}
Let $h=(i,j)$ be forced by a label $k$, whose occurrences in row $i$ and
column $j$ are $(i,c)$ and $(r,j)$.  Then $c<j$ if and only if $i<r$.
Moreover, all labels forcing the same hole have the same one of the two
orientations
\[
 c<j<{}\text{ and }i<r,\qquad\text{or}\qquad j<c\text{ and }r<i.
\]
\end{lemma}

\begin{proof}
The two occurrences of $k$ are comparable in the strict product order, so
$c<j$ holds exactly when $i<r$.  Now let $a<b$ force $h$.  Their positions in
row $i$ increase with the labels, and their positions in column $j$ do so as
well.  If $a$ were in the first orientation and $b$ in the second, their row
positions would have the correct order but their column positions would have
the reverse order.  If the orientations were reversed, the contradiction
would occur in the row.  Hence the orientation is common.
\end{proof}

Call the first orientation \emph{left-facing} and the second
\emph{right-facing}; let $\Hcal_L$ and $\Hcal_R$ be the corresponding sets of
forced holes.  For the ordered occurrences
$(r_1,c_1)<\cdots<(r_{t_k},c_{t_k})$ of label $k$, put
\[
 \ell^k_{pq}=(r_p,c_q),\qquad r^k_{pq}=(r_q,c_p)\qquad(p<q).
\]
The first cross hole lies in $\Hcal_L$ and the second in $\Hcal_R$.

\begin{definition}[Forcing graph]
The \emph{forcing graph} $\Gamma(A)$ is the edge-coloured bipartite graph with
parts $\Hcal_L$ and $\Hcal_R$.  For every label $k$ and $p<q$, it has an edge
of colour $k$ joining $\ell^k_{pq}$ to $r^k_{pq}$.
\end{definition}

A matching between subsets of two product-ordered grids is
\emph{order-preserving} if $x<x'$ for two first endpoints implies $y<y'$ for
their corresponding second endpoints.

\begin{theorem}[Forcing-graph normal form]\label{thm:forcing-normal-form}
For every monotone partial matrix $A$, the graph $\Gamma(A)$ has the following
properties.
\begin{enumerate}[label=(\alph*)]
\item It is simple and bipartite with parts $\Hcal_L,\Hcal_R$.
\item The degree of every forced hole $h$ is $d(h)$.
\item The neighbours of each vertex form a strict product-order chain; in that
chain order, the incident edge colours strictly increase.
\item Each colour class is an order-preserving matching.
\item Every four-cycle has four distinct edge colours.
\item Its edge and wedge counts are
\begin{equation}\label{eq:edge-wedge-counts}
 e(\Gamma(A))=\sum_{k=1}^n\binom{t_k}{2}
 =\frac12\sum_{h\in\Hcal(A)}d(h),
 \qquad
 \sum_{h\in V(\Gamma(A))}\binom{\deg h}{2}=E(A).
\end{equation}
\end{enumerate}
\end{theorem}

\begin{proof}
The graph is bipartite by construction.  Its two cross holes determine the
two opposite filled corners; those corners have the edge colour as their
common label.  Thus two colours cannot give the same edge, and the graph is
simple.

Fix a left-facing hole $h=(i,j)$.  A forcing label $k$ occurs at unique cells
$(i,c_k)$ and $(r_k,j)$ with $c_k<j$ and $i<r_k$.  The edge produced by these
two occurrences joins $h$ to $(r_k,c_k)$.  Conversely, every incident edge
arises in this way, and different labels give different neighbours.  Hence
$\deg h=d(h)$.  If $a<b$ both force $h$, row and column monotonicity give
$c_a<c_b$ and $r_a<r_b$.  Therefore the neighbours, and their incident
colours, increase together.  The right-facing case is identical after
reversing both inequalities.

For a fixed label $k$, each of the maps $(p,q)\mapsto\ell^k_{pq}$ and
$(p,q)\mapsto r^k_{pq}$ is injective, so its colour class is a matching.  If
$\ell^k_{pq}<\ell^k_{uv}$, then $p<u$ and $q<v$, and consequently
$r^k_{pq}<r^k_{uv}$.  The same argument with the order reversed proves that
the matching is order-preserving.

Consider a four-cycle.  Its two left vertices are comparable because both
belong to the neighbourhood chain of either right vertex; similarly, its two
right vertices are comparable.  Write them as $x<x'$ and $y<y'$.  If the
colours at $(x,y),(x,y'),(x',y),(x',y')$ are $a,b,c,d$, respectively, the
neighbourhood order gives
\[
 a<b<d,\qquad a<c<d.
\]
Adjacent edges have different colours because each colour class is a
matching.  The only possible repetition is $b=c$, but then that colour has
endpoints ordered as $x<x'$ on the left and $y'>y$ on the right, contrary to
the order-preserving property.  Thus the cycle is rainbow.

There is one edge for each unordered pair of occurrences of each label.  This
proves the first edge equality.  Lemma~\ref{lem:first-completion}, or the
degree identity and the handshake lemma, gives the second.  Finally, summing
$\binom{\deg h}{2}=\binom{d(h)}2$ over the forced holes gives the wedge
identity.
\end{proof}

The rainbow property extends from four-cycles to complete bipartite subgraphs.

\begin{corollary}[Rainbow bicliques]\label{cor:rainbow-bicliques}
Every $K_{s,t}$ contained in $\Gamma(A)$ is rainbow.  Consequently $st\le n$.
In particular, two vertices in the same part have at most $\lfloor n/2\rfloor$
common neighbours.
\end{corollary}

\begin{proof}
Adjacent edges have distinct colours.  Two disjoint edges of a biclique are
opposite edges of a four-cycle and hence also have distinct colours by
Theorem~\ref{thm:forcing-normal-form}(e).  Thus the $st$ edges have distinct
colours, while only $n$ labels are available.  The final assertion is the case
$s=2$.
\end{proof}

The next construction shows that no fixed biclique can be excluded.

\begin{proposition}[Arbitrary forcing bicliques]\label{prop:arbitrary-bicliques}
For every $s,t\ge2$, there is an order-$n$ monotone partial matrix with
$n=st$ whose forcing graph contains $K_{s,t}$.
\end{proposition}

\begin{proof}
For $a\in[s]$ and $b\in[t]$, set $\lambda_{ab}=a+s(b-1)$.  Fill two disjoint
rectangles by
\[
 A(a,b)=\lambda_{ab},\qquad
 A(s+b,t+a)=\lambda_{ab},
\]
and leave all other cells empty.  Both rectangles fit because $s+t\le st$.
Rows and columns increase, and the two occurrences of every label are strictly
ordered in both coordinates; hence the matrix is monotone.

For $a\in[s]$ put $x_a=(a,t+a)$, and for $b\in[t]$ put
$y_b=(s+b,b)$.  The two filled cells carrying $\lambda_{ab}$ are the other
corners of the rectangle with corners $x_a$ and $y_b$.  Therefore $x_ay_b$ is
an edge of colour $\lambda_{ab}$ for every $(a,b)$, giving the required
biclique.
\end{proof}

\section{Common fibres and moment inequalities}
\label{sec:common-fibres}

The forcing constraints also give exact inequalities for common fibres.  For
a set $I$ of rows, define
\begin{align*}
 C(I)&=\{j:\text{column }j\text{ is filled in every row of }I\},\\
 L(I)&=\{k:\text{label }k\text{ occurs in every row of }I\}.
\end{align*}
Let $r_i$ denote the number of filled cells in row $i$.

\begin{theorem}[Endpoint exclusion]\label{thm:endpoint-exclusion}
For every set $I$ of $p\ge2$ rows of an order-$n$ monotone partial matrix,
\begin{equation}\label{eq:endpoint-exclusion}
 p|L(I)|+|C(I)|\le n,\qquad
 |L(I)|+p|C(I)|\le n,\qquad
 |L(I)|+|C(I)|\le\min_{i\in I}r_i.
\end{equation}
The analogous statements obtained by permuting the three coordinate roles
also hold.
\end{theorem}

\begin{proof}
For $i\in I$ and $k\in L(I)$, let $c_i(k)$ be the column of the $k$-cell in
row $i$.  We claim that the $p|L(I)|$ values $c_i(k)$ are all distinct.
Suppose $c_i(k)=c_{i'}(\ell)=j$ for distinct pairs $(i,k)$ and $(i',\ell)$.
Equality is impossible if $i=i'$ or $k=\ell$, so assume $i<i'$.  Column
monotonicity gives $k<\ell$.  Since the $k$-cells form a strict product-order
chain, $c_i(k)<c_{i'}(k)$; row monotonicity in row $i'$ gives
$c_{i'}(k)<c_{i'}(\ell)$.  Hence $j<c_{i'}(k)<j$, a contradiction.

None of these column values lies in $C(I)$.  Indeed, suppose that the $k$-cell
in row $i$ lies in a common column $j$, and choose $i'\in I\setminus\{i\}$.
If $i<i'$, the entry at $(i',j)$ is larger than $k$, whereas the equal-label
chain places the $k$-cell of row $i'$ strictly to the right of $j$; row
monotonicity gives the opposite comparison.  If $i'<i$, both comparisons
reverse and again contradict each other.  Thus the $p|L(I)|$ values $c_i(k)$
and the $|C(I)|$ common columns are distinct, proving the first inequality.

Interchange the roles of column and label, retaining the row coordinate.
Proposition~\ref{prop:matrix-correspondence} gives another monotone partial
matrix, and the first inequality becomes $p|C(I)|+|L(I)|\le n$.  Finally, in
any row $i\in I$, the cells with labels in $L(I)$ and the cells in columns
$C(I)$ are disjoint by the preceding argument.  All are filled, proving the
third inequality.  Coordinate symmetry gives the remaining versions.
\end{proof}

Let $c_j$ and $t_k$ denote column and label degrees in the fixed
representation.

\begin{corollary}[Common-neighbour moments]\label{cor:common-neighbour-moments}
For every $2\le p\le n$,
\begin{align}
 p\sum_{k=1}^n\binom{t_k}{p}+\sum_{j=1}^n\binom{c_j}{p}
 &\le n\binom np,\label{eq:moment-one}\\
 \sum_{k=1}^n\binom{t_k}{p}+p\sum_{j=1}^n\binom{c_j}{p}
 &\le n\binom np,\label{eq:moment-two}\\
 \sum_{\substack{I\subseteq[n]\\|I|=p}}
 \bigl(|L(I)|+|C(I)|\bigr)
 &\le
 \sum_{\substack{I\subseteq[n]\\|I|=p}}\min_{i\in I}r_i.
 \label{eq:moment-three}
\end{align}
All cyclic variants hold.
\end{corollary}

\begin{proof}
A label of degree $t_k$ belongs to $L(I)$ for exactly $\binom{t_k}{p}$ choices
of $I$.  A column of degree $c_j$ belongs to $C(I)$ for exactly
$\binom{c_j}{p}$ choices.  Sum the three inequalities in
\eqref{eq:endpoint-exclusion} over all $p$-subsets $I$ of the rows.
\end{proof}

For $p=2$, if two rows share $D$ labels and $Q$ occupied columns, the first
two inequalities say $2D+Q\le n$ and $D+2Q\le n$.  These are independent of
the row degrees and strengthen the disjoint-cell bound
$D+Q\le\min(r_i,r_{i'})$ in complementary directions.

\section{Cyclic completion-energy packing}
\label{sec:cyclic-energy}

Return to a pairwise two-comparable set $M\subseteq[n]^3$ of size $m$.  For
$q\in\{x,y,z\}$, let $A^{(q)}$ be the cyclic representation with coordinate
$q$ as label, and define
\[
 E_q=E(A^{(q)})
 =\sum_{h\in\Hcal(A^{(q)})}\binom{d_q(h)}2.
\]

A \emph{strict inversion} in a monotone partial matrix is an unordered pair
of filled cells $(i,j)$ and $(r,c)$ with $i<r$, $j<c$, and with the label at
$(i,j)$ larger than the label at $(r,c)$.  It is \emph{saturated} if one of
its two cross cells is a hole forced by both endpoint labels.  Write
$\Inv(A)$ and $\Sat(A)$ for the sets of strict and saturated inversions.

\begin{lemma}[Saturated-inversion correspondence]\label{lem:saturated-inversion}
For every monotone partial matrix $A$, there is a bijection between
double-forcing events
\[
 (h,\{a,b\}),\qquad a<b,\quad a,b\in K_h,
\]
and saturated strict inversions.  Consequently
\begin{equation}\label{eq:saturated-inversion-count}
 E(A)=|\Sat(A)|\le |\Inv(A)|.
\end{equation}
\end{lemma}

\begin{proof}
Let $h=(i,j)$ be forced by $a<b$.  In the left-facing orientation, pair the
$b$-cell in row $i$ with the $a$-cell in column $j$.  The former is northwest
of the latter and has the larger label, so the pair is a strict inversion;
their other cross cell is $h$, forced by both labels.  In the right-facing
orientation, pair the $b$-cell in column $j$ with the $a$-cell in row $i$.
This again gives a saturated strict inversion.  This defines the forward map.

Conversely, a saturated inversion together with its saturated cross hole
recovers the two endpoint labels and hence the double-forcing event.  It
remains to prove that a strict inversion cannot have both cross holes
saturated.  Write its endpoints as a northwest cell $(p,q)$ of label $b$ and
a southeast cell $(r,s)$ of label $a$, where $p<r$, $q<s$, and $a<b$.
If the northeast cross hole $(p,s)$ is saturated, then row $p$ contains an
additional $a$-cell strictly left of $(p,q)$.  If the southwest cross hole
$(r,q)$ is saturated, then column $q$ contains an additional $a$-cell strictly
above $(p,q)$.  Those two additional $a$-cells are strictly anti-aligned,
contradicting the equal-$a$ chain condition.  Hence there is a unique
saturated cross hole, so the inverse is well defined.  The two constructions
are inverse to one another, proving bijectivity.
\end{proof}

\begin{theorem}[Cyclic energy packing]\label{thm:cyclic-energy-packing}
For every pairwise two-comparable $M\subseteq[n]^3$ of size $m$,
\begin{equation}\label{eq:cyclic-energy-packing}
 \boxed{E_x+E_y+E_z\le\binom m2.}
\end{equation}
More precisely, the saturated-inversion images in the three coordinate
representations are pairwise disjoint sets of unordered pairs from $M$.
\end{theorem}

\begin{proof}
Apply Lemma~\ref{lem:saturated-inversion} to the three representations.  A
strict inversion in $A^{(q)}$ is a pair of triples for which the two
non-$q$ coordinates change in the same direction and the $q$ coordinate
changes in the opposite direction.  Thus $q$ is the unique exceptional
coordinate of the pair.  The same unordered pair cannot have two distinct
unique exceptional coordinates, so the three saturated-inversion images are
disjoint.  Each image is contained in $\binom{M}{2}$, proving
\eqref{eq:cyclic-energy-packing}.
\end{proof}

By Theorem~\ref{thm:forcing-normal-form}, $E_q$ is the wedge count of the
forcing graph in representation $q$.  Thus the total number of wedges in the
three canonical graphs is packed into a single multiplicity-one resource:
the unordered pairs of members of $M$.

\section{Exact slack and consequences}
\label{sec:exact-slack}

Let $t_u^{(q)}$ be the number of members of $M$ whose $q$ coordinate is $u$,
and put
\begin{equation}\label{eq:def-Sq}
 S_q=\sum_{u=1}^n t_u^{(q)}(t_u^{(q)}-1).
\end{equation}
Every representation has $H=n^2-m$ holes.  For $n\ge2$ one has $H>0$:
otherwise Lemma~\ref{lem:first-completion} would force every label degree to be
at most one, contrary to $m=n^2>n$.  Define
\begin{equation}\label{eq:def-variances}
 V_q=\sum_{u=1}^n\left(t_u^{(q)}-\frac mn\right)^2,\qquad
 W_q=\sum_{h\in\Hcal(A^{(q)})}
 \left(d_q(h)-\frac{S_q}{H}\right)^2.
\end{equation}
Let $P_{\parallel}$ count unordered pairs of members of $M$ whose three
coordinates are distinct and all change in the same direction.  Finally, set
\[
 U_q=|\Inv(A^{(q)})|-E_q,
\]
the number of unsaturated strict inversions in the $q$ representation.

\begin{theorem}[Exact slack identity]\label{thm:exact-slack}
For every pairwise two-comparable $M\subseteq[n]^3$ of size $m$ and $n\ge2$,
\begin{equation}\label{eq:exact-slack}
 \boxed{
 \binom m2
 =\frac1{2H}\sum_{q\in\{x,y,z\}}S_q^2
 +\frac12\sum_{q\in\{x,y,z\}}W_q
 +P_{\parallel}+\sum_{q\in\{x,y,z\}}U_q.}
\end{equation}
Moreover,
\begin{equation}\label{eq:coordinate-slack}
 S_q=\frac{m^2}{n}-m+V_q.
\end{equation}
All remainder terms $V_q,W_q,P_{\parallel},U_q$ are nonnegative.
\end{theorem}

\begin{proof}
Lemma~\ref{lem:first-completion} gives
$\sum_{h\in\Hcal(A^{(q)})}d_q(h)=S_q$.  Expanding $W_q$ therefore yields the
exact variance identity
\begin{equation}\label{eq:energy-variance}
 E_q=\frac12\left(\frac{S_q^2}{H}+W_q-S_q\right).
\end{equation}

We next partition the unordered pairs from $M$.  Two distinct members cannot
agree in two coordinates, so the number of pairs agreeing in one coordinate
is counted exactly once by the coordinate fibres and equals
\[
 \sum_q\sum_{u=1}^n\binom{t_u^{(q)}}2=\frac12\sum_q S_q.
\]
A pair differing in all three coordinates is either counted by
$P_{\parallel}$ or has a unique exceptional coordinate $q$, in which case it
is a strict inversion in $A^{(q)}$.  Hence
\begin{equation}\label{eq:pair-partition}
 \binom m2=\frac12\sum_qS_q+P_{\parallel}
 +\sum_q(E_q+U_q).
\end{equation}
Substituting \eqref{eq:energy-variance} cancels the linear $S_q$ terms and
proves \eqref{eq:exact-slack}.  Finally,
\[
 V_q=\sum_u(t_u^{(q)})^2-\frac{m^2}{n},
\]
and subtracting $\sum_u t_u^{(q)}=m$ gives
\eqref{eq:coordinate-slack}.  Nonnegativity follows from the definitions and
Lemma~\ref{lem:saturated-inversion}.
\end{proof}

The identity keeps both variance defects and the full unused part of cyclic
inversion packing.  In particular,
\begin{equation}\label{eq:basic-energy-lower}
 S_q\ge \frac{m^2}{n}-m,\qquad
 E_q\ge\frac12\left(\frac{S_q^2}{H}-S_q\right).
\end{equation}

\subsection{Density calibration and stability}

\begin{corollary}[Elementary density calibration]\label{cor:density-calibration}
One has
\[
 \limsup_{n\to\infty}\frac{G(n)}{n^2}
 \le \rho_0:=\frac{\sqrt{13}-1}{6}=0.434258\ldots.
\]
\end{corollary}

\begin{proof}
Consider a sequence with $m/n^2\to\rho>0$.  The case $\rho=1$ is impossible:
then $H=o(n^2)$ and $S_q\ge m^2/n-m=(1-o(1))n^3$, so
\eqref{eq:basic-energy-lower} would give $E_q/n^4\to\infty$, contrary to
Theorem~\ref{thm:cyclic-energy-packing}.  Thus $0<\rho<1$.  For every $q$,
\[
 S_q\ge B:=\frac{m^2}{n}-m=\Theta(n^3),\qquad H=\Theta(n^2).
\]
The function $s\mapsto s^2/H-s$ is increasing on $[B,\infty)$ for all
sufficiently large $n$.  It follows from \eqref{eq:basic-energy-lower} that
\[
 \frac{E_q}{n^4}\ge\frac{\rho^4}{2(1-\rho)}+o(1).
\]
On the other hand, Theorem~\ref{thm:cyclic-energy-packing} gives
$\sum_qE_q/n^4\le\rho^2/2+o(1)$.  Therefore
\[
 \frac{3\rho^4}{1-\rho}\le\rho^2,
 \qquad\text{equivalently}\qquad 3\rho^2+\rho-1\le0.
\]
The positive root is $\rho_0$.
\end{proof}

The known bound $G(n)=o(n^2)$ is asymptotically stronger.  The corollary is
included only as a calibration of the cyclic identity, not as a competitive
upper bound.

\begin{corollary}[Slack stability]\label{cor:slack-stability}
If a sequence of pairwise two-comparable sets satisfies
$m/n^2\to\rho_0$, then, for every $q\in\{x,y,z\}$,
\[
 V_q=o(n^3),\qquad W_q=o(n^4),\qquad
 P_{\parallel}+\sum_qU_q=o(n^4).
\]
\end{corollary}

\begin{proof}
Put $B=m^2/n-m$, so $S_q=B+V_q$.  Subtract $3B^2/(2H)$ from both sides of
\eqref{eq:exact-slack}.  The remainder on the right is the sum of the
nonnegative quantities
\[
 \frac1{2H}\sum_q(2BV_q+V_q^2)
 +\frac12\sum_qW_q+P_{\parallel}+\sum_qU_q.
\]
Since $3\rho_0^2+\rho_0-1=0$, the left side is $o(n^4)$.  Also
$B=\Theta(n^3)$ and $H=\Theta(n^2)$, so each asserted estimate follows.
\end{proof}

In view of the known $o(n^2)$ upper bound, the hypothesis of
Corollary~\ref{cor:slack-stability} does not describe an extremal sequence.
Its role is to record the equality conditions of the elementary cyclic
argument.

\subsection{A low-energy criterion}

\begin{corollary}[Low-energy power saving]\label{cor:low-energy}
Fix $\eta>0$.  If $E_q=O(n^{4-\eta})$ for at least one coordinate $q$, then
\begin{equation}\label{eq:low-energy-bound}
 m=O\bigl(n^{2-\eta/4}+n^{3/2}\bigr).
\end{equation}
Consequently, if $m=n^{2-o(1)}$, then in every cyclic representation
\[
 e(\Gamma(A^{(q)}))=n^{3-o(1)},
 \qquad
 \sum_h\binom{\deg_{\Gamma(A^{(q)})}h}{2}=n^{4-o(1)}.
\]
\end{corollary}

\begin{proof}
If $m<2n^{3/2}$, then \eqref{eq:low-energy-bound} is immediate.  Otherwise
$S_q\ge m^2/(2n)\ge2n^2$.  Since $H\le n^2$, we have
$S_q^2/H\ge2S_q$, and \eqref{eq:basic-energy-lower} gives
\[
 E_q\ge\frac{S_q^2}{4H}
 \ge\frac{S_q^2}{4n^2}
 \ge\frac{m^4}{16n^4}.
\]
Comparison with $E_q=O(n^{4-\eta})$ gives
$m=O(n^{2-\eta/4})$.  For the final assertion,
\eqref{eq:first-completion} and \eqref{eq:coordinate-slack} give
$e(\Gamma(A^{(q)}))=S_q/2=n^{3-o(1)}$, while the preceding lower bound and
Theorem~\ref{thm:forcing-normal-form}(f) give the wedge estimate.
\end{proof}

\section{Recursive pressure overlap}
\label{sec:pressure}

Let a forced hole $h=(i,j)$ have forcing labels
$K_h=\{k_1<\cdots<k_d\}$.  Write
\[
 A(i,c_p)=k_p=A(r_p,j),
 \qquad
 P_h=\{(r_p,c_q):1\le p,q\le d\}.
\]
We call $P_h$ the \emph{pressure grid} of $h$.  By
Lemma~\ref{lem:common-orientation}, either
\begin{equation}\label{eq:left-orientation}
 c_1<\cdots<c_d<j,\qquad i<r_1<\cdots<r_d,
\end{equation}
or
\begin{equation}\label{eq:right-orientation}
 j<c_1<\cdots<c_d,\qquad r_1<\cdots<r_d<i.
\end{equation}

\begin{lemma}[Triangular pressure bound]\label{lem:triangular-pressure}
Every forced hole satisfies
\begin{equation}\label{eq:triangular-pressure}
 |P_h\cap\Hcal(A)|\ge\frac{d(h)(d(h)+1)}2.
\end{equation}
\end{lemma}

\begin{proof}
In the orientation \eqref{eq:left-orientation}, every cell $(r_p,c_q)$ with
$p\le q$ is a hole.  Indeed, a label placed there would be smaller than $k_p$
by row monotonicity and larger than $k_q$ by column monotonicity, which is
impossible because $k_p\le k_q$.  In the orientation
\eqref{eq:right-orientation}, the same argument, with inequalities reversed,
shows that all $(r_p,c_q)$ with $p\ge q$ are holes.  Either triangular half
has $d(d+1)/2$ cells.
\end{proof}

Fix a physical cell $z=(r,c)$, whether filled or empty.  Define the signed
sets of pressure centres
\begin{align}
 \mathcal Z_z^+&=\{(i,j)\in\Hcal(A):z\in P_{(i,j)},\ i<r,\ c<j\},\label{eq:Zplus}\\
 \mathcal Z_z^-&=\{(i,j)\in\Hcal(A):z\in P_{(i,j)},\ r<i,\ j<c\}.
 \label{eq:Zminus}
\end{align}
These are the only possibilities, corresponding to
\eqref{eq:left-orientation} and \eqref{eq:right-orientation}.  Put
\begin{equation}\label{eq:def-loads}
 L_0(z)=|\mathcal Z_z^+|+|\mathcal Z_z^-|,
 \qquad
 L_1(z)=\sum_{h:z\in P_h}d(h).
\end{equation}

We use the following finite form of the Loomis--Whitney inequality.

\begin{lemma}[Projection inequality]\label{lem:projection-inequality}
If $X$ is a finite subset of a three-fold Cartesian product, then
\[
 |X|^2\le |\pi_{12}X|\,|\pi_{13}X|\,|\pi_{23}X|.
\]
\end{lemma}

\begin{proof}
For $(a,b)\in\pi_{12}X$, let $X_{ab}=\{c:(a,b,c)\in X\}$.  Cauchy--Schwarz
gives
\[
 |X|^2\le |\pi_{12}X|\sum_{a,b}|X_{ab}|^2.
\]
The sum counts ordered pairs $(a,b,c,c')$ with $(a,b,c),(a,b,c')\in X$.
The map
\[
 (a,b,c,c')\longmapsto((a,c),(b,c'))
\]
is injective into $\pi_{13}X\times\pi_{23}X$.  Hence the sum is at most
$|\pi_{13}X|\,|\pi_{23}X|$, proving the result.
\end{proof}

\begin{theorem}[Recursive pressure support]\label{thm:recursive-pressure}
For every cell $z$ and each sign $\varepsilon\in\{+,-\}$, the centre set
$\mathcal Z_z^\varepsilon$ has two canonical injective encodings as pairwise
two-comparable subsets of $[n]^3$.  If $A$ has $m$ filled cells, then
\begin{equation}\label{eq:pointwise-pressure}
 \boxed{L_0(z)\le2m,\qquad L_1(z)\le2m^{3/2}.}
\end{equation}
More precisely, each signed contribution to $L_1(z)$ is at most $m^{3/2}$.
\end{theorem}

\begin{proof}
We prove the assertions for $\mathcal Z_z^+$; reversing all inequalities gives
the proof for $\mathcal Z_z^-$.  Let $h=(i,j)\in\mathcal Z_z^+$.  Since
$z=(r,c)\in P_h$, the boundary cells
\[
 A(r,j)=a_j,\qquad A(i,c)=b_i
\]
are filled and both labels belong to $K_h$.  For fixed $z$, the first label
depends only on $j$ and increases strictly with $j$; the second depends only
on $i$ and increases strictly with $i$.

Because $a_j$ forces $h$ and already occurs in column $j$, it has a unique
occurrence $(i,p_{ij})$ in row $i$.  Similarly, $b_i$ has a unique occurrence
$(q_{ij},j)$ in column $j$.  Consider the two encodings
\begin{equation}\label{eq:pressure-encodings}
 (i,j)\longmapsto(i,j,p_{ij}),
 \qquad
 (i,j)\longmapsto(i,j,q_{ij}).
\end{equation}
We verify the first.  If two images have the same first coordinate $i$, then
$j<j'$ implies $a_j<a_{j'}$ and row monotonicity gives
$p_{ij}<p_{ij'}$.  If they have the same second coordinate $j$, the
equal-$a_j$ chain gives $i<i'\Rightarrow p_{ij}<p_{i'j}$.  If they have the
same third coordinate $p$, then column monotonicity in column $p$ gives
$i<i'\Rightarrow a_j<a_{j'}$, and row monotonicity in row $r$ gives $j<j'$.
If all coordinates differ, one of the two triples is smaller in at least two
coordinates automatically.  Hence the image is pairwise two-comparable.  The
second encoding follows in the same way, using successively the equal-$b_i$
chain, column monotonicity, and row monotonicity.

The first encoding assigns to $(i,j)$ the filled cell $(i,p_{ij})$.  These
cells are distinct: centres in different rows have different first
coordinates, while for fixed $i$ the preceding argument gives distinct
$p_{ij}$.  Thus $|\mathcal Z_z^+|\le m$.  The other sign gives
$L_0(z)\le2m$.

For the weighted assertion, form the incidence set
\[
 \mathcal I_z^+=\{(i,j,k):(i,j)\in\mathcal Z_z^+,\ k\in K_{(i,j)}\}.
\]
Its $(i,j)$-projection is the centre set and has size at most $m$.  Its
$(i,k)$-projection is contained in the set of occupied row--label pairs of
$A$, and its $(j,k)$-projection is contained in the occupied column--label
pairs.  Each of the latter sets has size exactly $m$, since a label occurs at
most once in a fixed row or column.  Lemma~\ref{lem:projection-inequality}
therefore gives
\[
 |\mathcal I_z^+|^2
 \le |\pi_{ij}\mathcal I_z^+|\,
      |\pi_{ik}\mathcal I_z^+|\,
      |\pi_{jk}\mathcal I_z^+|
 \le m^3.
\]
Since $|\mathcal I_z^+|=\sum_{h\in\mathcal Z_z^+}d(h)$, the positive signed
load is at most $m^{3/2}$.  The negative signed load obeys the same bound;
adding them completes the proof.
\end{proof}

The theorem turns uncontrolled geometric overlap into a recursive incidence
problem on genuine two-comparable supports.  It also yields an explicit
sufficient condition for a power saving.

\begin{corollary}[Pressure-load criterion]\label{cor:pressure-load}
Let
\[
 \Omega_1(A)=\max_{z\in\Hcal(A)}L_1(z),\qquad
 \Omega_0(A)=\max_{z\in\Hcal(A)}L_0(z).
\]
Suppose a class of order-$n$ monotone partial matrices satisfies, uniformly,
$\Omega_1(A)\le Cn^{3-\eta}$ for constants $C<\infty$ and $\eta>0$.  Then
every member of the class has
\begin{equation}\label{eq:pressure-power-saving}
 m=O_C(n^{2-\eta/6}+n).
\end{equation}
Under the alternative hypothesis $\Omega_0(A)\le Cn^{2-\eta}$, one has
$m=O_C(n^{2-\eta/4}+n)$.
\end{corollary}

\begin{proof}
Write $S=\sum_{h\in\Hcal(A)}d(h)=\sum_k t_k(t_k-1)$ and $H=n^2-m$.
Double-count the weighted incidences $(h,z)$ with
$z\in P_h\cap\Hcal(A)$.  Lemma~\ref{lem:triangular-pressure} and the
power-mean inequality give
\[
 H\Omega_1(A)
 \ge\frac12\sum_h\bigl(d(h)^3+d(h)^2\bigr)
 \ge\frac{S^3}{2H^2}.
\]
If $m\ge2n$, then
$S\ge m^2/n-m\ge m^2/(2n)$, and $H\le n^2$.  Consequently
\[
 m^6\le16n^9\Omega_1(A)\le16C n^{12-\eta},
\]
so $m\le(16C)^{1/6}n^{2-\eta/6}$.  The case $m<2n$ is absorbed by the
$O(n)$ term.

For the unweighted count, the same double counting gives
\[
 H\Omega_0(A)\ge\frac12\sum_h(d(h)^2+d(h))\ge\frac{S^2}{2H}.
\]
When $m\ge2n$, this yields
$m^4\le8n^6\Omega_0(A)\le8C n^{8-\eta}$ and hence the asserted exponent
$2-\eta/4$.
\end{proof}

The exponent $3/2$ in Theorem~\ref{thm:recursive-pressure} is optimal as a
function of $m$.

\begin{proposition}[Sharp pressure overlap]\label{prop:sharp-pressure}
For all integers $s,t\ge1$, put
$N=s(s+t)+s+1$.  There is an order-$N$ monotone partial matrix with
\[
 m=2s(s+t+1)
\]
filled cells and a hole $z$ such that
\[
 L_0(z)\ge s^2,\qquad L_1(z)\ge s^2(t+2).
\]
Taking $t=s$ gives $m=\Theta(N)$ and
$L_1(z)=\Theta(N^{3/2})=\Theta(m^{3/2})$.
\end{proposition}

\begin{proof}
Use rows and columns in $[N]$.  Set $c_0=1$, $r_0=N$, and
\[
 C_{i,u}=1+(i-1)(s+t)+u,\qquad
 J_j=1+s(s+t)+j,\qquad
 R_{j,u}=s+(j-1)(s+t)+u.
\]
For $i,j\in[s]$ and $u\in[s+t]$, let
\[
 B_i=i,\qquad K_u=s+u,\qquad A_j=s+t+j=K_{t+j}.
\]
Fill
\begin{align*}
 A(i,c_0)&=B_i,& A(i,C_{i,u})&=K_u,\\
 A(R_{j,u},J_j)&=u,& A(r_0,J_j)&=A_j,
\end{align*}
and leave every other cell empty.  The listed indices lie in $[N]$.  Each
nonempty row and column is strictly increasing.  It remains to check the
equal-label chains.  A label $u\in[s]$ occurs first at $(u,c_0)$ and then at
the cells $(R_{j,u},J_j)$, which move strictly southeast with $j$.  A label
$s+u$ with $u\in[t]$ occurs at $(i,C_{i,u})$ for $i\in[s]$ and then at
$(R_{j,s+u},J_j)$ for $j\in[s]$; both subfamilies are southeast chains, and
every cell in the first is northwest of every cell in the second.  Finally, a
label $s+u$ with $t<u\le s+t$ occurs at $(i,C_{i,u})$ for $i\in[s]$ and once
more at $(r_0,J_{u-t})$.  These cells also form a southeast chain.  Hence $A$
is monotone.

For every $(i,j)\in[s]^2$, the cell $(i,J_j)$ is a hole forced by exactly
\[
 \{B_i,A_j,K_1,\ldots,K_t\}.
\]
Its pressure grid contains $z=(r_0,c_0)$, using the $A_j$ column arm and the
$B_i$ row arm.  Therefore these $s^2$ centres contribute $s^2$ to $L_0(z)$
and $s^2(t+2)$ to $L_1(z)$.  The two displayed families of filled cells are
disjoint, each having $s(s+t+1)$ members, which proves the formula for $m$.
Putting $t=s$ gives the stated asymptotics.
\end{proof}

\section{Scope and the multiscale obstruction}
\label{sec:scope}

The results above do not prove \eqref{eq:open-problem}.  They instead isolate
three limitations of one-scale arguments.  First, forcing graphs need not be
acyclic.  For example,
\[
 \begin{pmatrix}
 1&3&\cdot&\cdot\\
 2&4&\cdot&\cdot\\
 \cdot&\cdot&1&2\\
 \cdot&\cdot&3&4
 \end{pmatrix}
\]
is monotone and its forcing graph contains a four-cycle; the four colours are
distinct, as Theorem~\ref{thm:forcing-normal-form} requires.  More generally,
Proposition~\ref{prop:arbitrary-bicliques} excludes every strategy based only
on forbidding one fixed biclique.

Second, Proposition~\ref{prop:sharp-pressure} shows that the exponent in the
unconditional bound $L_1(z)\le2m^{3/2}$ cannot be improved using only $m$.
When $m$ is near the quadratic scale, this bound is of cubic order, exactly
the scale at which Corollary~\ref{cor:pressure-load} ceases to save a power.
Any improvement must therefore use the geometry or density of the recursive
centre encodings, not only their cardinality.

Third, the baseline part of the exact slack identity is too small at a
slowly subquadratic scale.  If formally $m=n^2/L$ with $L\to\infty$ and
$L=o(n)$, then $H=(1-o(1))n^2$ and
\[
 \frac{1}{2H}\sum_q\left(\frac{m^2}{n}-m\right)^2
 =\Theta\!\left(\frac{n^4}{L^4}\right),
 \qquad
 \binom m2=\Theta\!\left(\frac{n^4}{L^2}\right).
\]
Thus the exact identity can exclude this scale only if it forces one of the
variance, alignment, or unsaturated-inversion terms to be large in a way that
persists under restriction.  No such cross-scale estimate is proved here.

The recursive pressure theorem supplies a natural carrier for such an
estimate: grids through one cell already produce smaller two-comparable sets.
A fixed power saving would follow from a density decrement on those supports,
from a multiplicative loss on a positive proportion of scales, or from a new
interaction among the three cyclic representations.  Establishing any one of
these remains open.

\appendix

\section{Finite verification}
\label{app:verification}

The accompanying program \texttt{src/verify\_finite.py} is a supplementary
consistency check; it is not used in any proof.  It enumerates the
$4^9=262144$ arrays in
\[
 \{\text{empty},1,2,3\}^{[3]\times[3]},
\]
retains exactly the $712$ arrays satisfying the three monotonicity conditions,
and checks the following statements for every retained array and every cyclic
coordinate representation:
\begin{enumerate}[label=(\roman*)]
\item common orientation, forcing-graph simplicity, degree equality, chain
neighbourhoods, order-preserving colour matchings, rainbow four-cycles, and
the edge and wedge identities;
\item endpoint exclusion for every row subset of size at least two;
\item the saturated-inversion correspondence, cyclic energy packing, the
pair partition \eqref{eq:pair-partition}, and the exact rational identity
\eqref{eq:exact-slack};
\item the two pressure-support encodings and all three projection bounds.
\end{enumerate}
The same program verifies Proposition~\ref{prop:arbitrary-bicliques} for four
small parameter pairs and Proposition~\ref{prop:sharp-pressure} for four
small parameter pairs.  All variance calculations use exact rational
arithmetic.  The complete command and expected output are given in the
repository README.  The finite computation proves only the stated finite
enumeration and checks the implementation against the general proofs; the
theorems themselves are not computer-assisted.

\begin{thebibliography}{99}

\bibitem{AronovEtAl}
B. Aronov, V. Dujmovi\'c, P. Morin, A. Ooms, and L. F. S. X. da Silveira,
More Turán-type theorems for triangles in convex point sets,
\emph{Electron. J. Combin.} \textbf{26} (2019), Paper No. 1.8, 26 pp.

\bibitem{Fox}
J. Fox,
A new proof of the graph removal lemma,
\emph{Ann. of Math. (2)} \textbf{174} (2011), 561--579.

\bibitem{FoxSudakovWigderson}
J. Fox, B. Sudakov, and Y. Wigderson,
Color-avoiding Ramsey for directed paths and vector sequences,
preprint, arXiv:2512.10438, 2025; revised 2026.

\bibitem{GowersLong}
W. T. Gowers and J. Long,
The length of an $s$-increasing sequence of $r$-tuples,
\emph{Combin. Probab. Comput.} \textbf{30} (2021), 686--721.

\bibitem{Loh}
P.-S. Loh,
Directed paths: from Ramsey to Ruzsa and Szemerédi,
preprint, arXiv:1505.07312, 2015; revised 2016.

\bibitem{OstergardPollanen}
P. R. J. Östergård and A. Pöllänen,
New results on tripod packings,
\emph{Discrete Comput. Geom.} \textbf{61} (2019), 271--284.

\bibitem{Stein}
S. K. Stein,
Packing tripods,
\emph{Math. Intelligencer} \textbf{17} (1995), no.~2, 37--39.

\bibitem{Tiskin}
A. Tiskin,
Packing tripods: narrowing the density gap,
\emph{Discrete Math.} \textbf{307} (2007), 1973--1981.

\end{thebibliography}

\end{document}
